\section{Introdução}
\label{sec:introducao}

O Sistema Integrado de Gestão de Atividades Acadêmicas (SIGAA), desenvolvido pela Universidade Federal do Rio Grande do Norte (UFRN) e adotado por diversas universidades federais, é o sistema pelo qual o estudante de graduação da Universidade de Brasília (UnB) consulta notas, horários, histórico e índices acadêmicos, emite documentos, realiza matrícula e acessa as turmas virtuais \cite{sigaa}. Não há alternativa institucional: uma dificuldade no sistema se converte em atraso, erro de planejamento ou dependência de colegas e de ferramentas de terceiros.

Há indícios de que o SIGAA não atende bem esse público. A interface foi concebida para computador de mesa, com largura fixa e menus que se abrem pela passagem do ponteiro do mouse; o horário das turmas aparece em notação codificada (``24M34''); e estudantes recorrem a ferramentas criadas pela comunidade para planejar o fluxo de disciplinas. Esses indícios, porém, são percepções. Para apoiar decisões de melhoria, precisam ser transformados em evidências mensuráveis e comparadas a critérios definidos antes da coleta.

A família SQuaRE oferece o arcabouço para essa transformação: um modelo de qualidade \cite{iso25010_2023}, um catálogo de medidas \cite{iso25023} e um processo de avaliação \cite{iso25040}. O método \textit{Goal-Question-Metric} (GQM) liga os objetivos às questões e às métricas \cite{basili1994}, e o \textit{Practical Software and Systems Measurement} (PSM) orienta a medição do próprio projeto \cite{mcgarry2002}.

O objetivo geral é \textbf{avaliar a qualidade do portal do discente de graduação do SIGAA/UnB, do ponto de vista do estudante}, para identificar com evidências os pontos fracos nas tarefas mais frequentes e apoiar a priorização de melhorias pela Secretaria de Tecnologia da Informação (STI/UnB). Os objetivos específicos são: (i) estabelecer os requisitos da avaliação (Fase 1); (ii) especificá-la pelo GQM, com métricas, níveis de pontuação e critérios de julgamento (Fase 2); (iii) projetá-la e executá-la, julgando os resultados e as relações entre as características (Fases 3 e 4); e (iv) executar uma ação de melhoria baseada nas evidências. Este artigo corresponde à Entrega 1 e cobre os itens (i) e (ii) \cite{ramos2026roteiro}.

A Seção~\ref{sec:referencial} apresenta o referencial teórico; a Seção~\ref{sec:metodologia}, a metodologia e a gestão do projeto; as Seções~\ref{sec:fase1} e~\ref{sec:fase2}, as Fases 1 e 2; e a Seção~\ref{sec:consideracoes}, as considerações parciais. A contribuição da equipe está no Apêndice~\ref{ap:contribuicao}. Os artefatos citados (páginas detalhadas, dados, instrumentos e histórico de versões) estão no repositório da equipe\footnote{\repositorio} e na página do projeto\footnote{\gitpage}.
